\cotefile{37}
\scanpage{3}
SGA : Manuscripts à frapper

\diagram{\begingroup\small
\(\begin{array}{llll}
\text{\deleted{SGA 1}} & \text{\deleted{Index notations et terminologique:}} & & \text{\deleted{\uline{10} p}}\\
\text{SGA \unsure{\deleted{2}4|\deleted{4}2}} & \text{\tikz[remember picture,baseline=(sgaiv.base)]\node[inner sep=0pt] (sgaiv) {Exp IV (fin)};} & \text{Verdier} & \text{\uline{\unsure{5|\deleted{3}5}0}}\\
& \text{Exp V} & \text{Verdier} & 80\smash{\raisebox{-1.5em}{\handwritten{\(\left.\rule{0pt}{3.6em}\right\}\)}}}\smash{\raisebox{-3em}{\handwritten{\(\left.\rule{0pt}{6.6em}\right]\)}}}\\
& \text{Exp VI (fin)} & \text{Verdier} & 50\\
& \text{Exp VI B} & \text{S\deleted{q}aint Donat} & \text{\unsure{\deleted{5}1|\deleted{1}5}0}\\
& \text{\shortstack[l]{Tables des matières pour les trois\\ volumes, index}} & & 30\\
& \text{Appendice à \unsure{XVII|\deleted{I}XVII}} & \text{Saint Donat} & 10\\
\text{SGA 5} & \text{Exp IV} & ? & 100\ ?\\
& \text{\tikz[remember picture,baseline=(sgavii.base)]\node[inner sep=0pt] (sgavii) {Exp VII};} & \text{Jouanolou} & \text{\uline{50}}\\
& \text{\tikz[remember picture,baseline=(sgaxi.base)]\node[inner sep=0pt] (sgaxi) {Exp XI};} & \text{Bucur} & \text{\uline{70}}\\
& \text{Exp XII} & \text{Bucur} & 100\smash{\raisebox{-0.65em}{\handwritten{\(\left.\rule{0pt}{2em}\right]\)}}}\\
& \text{Exp XIV} & \text{Grothendieck} & 20\\
& \text{\shortstack[l]{Introduction, table des matières,\\ index}} & & 15\\
\text{\deleted{SGA 6}} & \text{\deleted{Exp XI}} & \text{\deleted{Ferrand}} & \text{\deleted{70 ?}}\\
& \text{\deleted{Table des matières, index}} & & \text{\deleted{10}}\\
\text{SGA 7} & \text{Exp I à \deleted{II}IV} & ? & 120\ ?\\
& \text{Exp \deleted{V}V} & \text{Mme Raynaud} & \text{\tikz[remember picture,baseline=(raynaudcount.base)]\node[inner sep=2pt] (raynaudcount) {30};}\\
& \text{\tikz[remember picture,baseline=(sgarim.base)]\node[inner sep=0pt] (sgarim) {Exp VI};} & \text{Rim} & \text{\tikz[remember picture,baseline=(rimcount.base)]\node[inner sep=2pt] (rimcount) {\uline{70}};}\\
& \text{\shortstack[l]{Introduction, table des matières,\\ index}} & & 15
\end{array}\)
\endgroup}
% Right margin, opposite the cancelled SGA 1 row.
\margin{\handwritten{--- \unsure{tel|Cl} Saarbach}}
% Right margin, opposite the bracket joining SGA 5 Exp XII and Exp XIV.
\margin{\handwritten{--- \uline{écrire Bucur}}}
% Right margin, opposite the framed SGA 7 counts 30 and 70.
\margin{\handwritten{\uline{écr. Deligne}}}
% Pencil strokes are anchored to their source table rows; the counts remain typed.
\handwritten{\figure{\begin{tikzpicture}[remember picture,overlay]
\draw ([xshift=-.8cm,yshift=-.45cm]sgaiv.south west) -- ([yshift=-.08cm]sgaiv.south west) -- ([yshift=-.08cm]sgaiv.south east);
\draw ([xshift=-1.05cm,yshift=-.3cm]sgavii.west) -- (sgavii.west);
\draw ([xshift=-.95cm,yshift=-.4cm]sgaxi.west) -- (sgaxi.west);
\draw ([xshift=-1.4cm,yshift=-.55cm]sgarim.west) -- (sgarim.west);
\draw (raynaudcount.north west) -- (raynaudcount.north east) -- (rimcount.south east) -- (rimcount.south west) -- cycle;
\draw ([xshift=-1pt,yshift=1pt]raynaudcount.north west) -- ([xshift=1pt,yshift=1pt]raynaudcount.north east) -- ([xshift=1pt,yshift=-1pt]rimcount.south east) -- ([xshift=-1pt,yshift=-1pt]rimcount.south west) -- cycle;
\draw ([xshift=-3pt]raynaudcount.south west) -- ([xshift=1pt,yshift=-1pt]raynaudcount.south east);
\end{tikzpicture}}}

NB Les chiffres soulignés correspondent à des exposés dont le manuscrit est disponible dès maintenant; ceux suivis d'un signe ? à ceux dont il n'est pas certain qu'ils seront rédigés. En plus de la frappe d'exposés, tables de matières, index, il y a pour chacun des cinq séminaires envisagés un certain travail de mise\deleted{\unsure{r}} au point (\unsure{correction|corection} de fautes de frappe et autres erreurs, insertion de notes de bas de page etc), représentant un travail variable d'un séminaire à l'autre. Il faudrait sans doute fixer un forfait pour chacun d'eux, pour la rénumération de la sécrétaire, après examen avec elle des corrections à faire ...

Il y a actuellement dans les 250 pages prêtes pour la frappe. J'espère que dans les deux ou trois mois qui v\unsure{\deleted{ei}ie|\deleted{ie}ei}nnent, il y aura aussi les env. \unsure{\deleted{3}2|\deleted{2}3}00 pages qui manquent encore pour SGA 4 (qui sont déjà partiellement rédigés ...), ce qui ferait dans les 450 à 500 pages à frapper pour les cinq à six mois qui viennent. Cependant je ne peux le garantir, cela dépendra de Mr. Verdier, qui est à Brandeis cette année. Je me mets en rapport avec lui ces jours-ci.
% archive sheet number: 2
